\documentclass{article}

\usepackage{amsmath, amsthm, amssymb}
\usepackage{graphicx}
\usepackage{verbatim}
\usepackage{natbib}
\usepackage{caption}
\usepackage{subcaption}
\usepackage{fancyvrb}
\usepackage{enumerate}
\usepackage{relsize}
\usepackage{multirow}
\usepackage{mathrsfs}

\usepackage{hyperref}
\usepackage[margin=1.4in]{geometry}
\hypersetup{colorlinks,citecolor=blue,urlcolor=blue,linkcolor=blue}
\newcommand{\Lihua}[1]{{\color{blue}#1}}

\newcommand\blfootnote[1]{%
  \begingroup
  \renewcommand\thefootnote{}\footnote{#1}%
  \addtocounter{footnote}{-1}%
  \endgroup
}

\usepackage{stefan_tex}
\graphicspath{{./figures/}}

%\numberwithin{equation}{section}

%%%%%%%%% Theorems

\theoremstyle{definition}
\newtheorem{exam}{Example}
\newtheorem{defi}{Definition}
\newtheorem{assumption}{Assumption}
\newtheorem{property}{Property}

 \theoremstyle{plain}
 \newtheorem{prop}{Proposition}
 \newtheorem{conj}[prop]{Conjecture}
 \newtheorem{coro}[prop]{Corollary}
 \newtheorem{lemm}[prop]{Lemma}
 \newtheorem{theo}[prop]{Theorem}
%\theoremstyle{plain}
%\newtheorem{prop}{Proposition}[section]
%\newtheorem{coro}{Corollary}[section]
%\newtheorem{lemm}{Lemma}[section]
%\newtheorem{theo}{Theorem}[section]

\theoremstyle{remark}
\newtheorem{comm}{Comment}
\newtheorem{rema}{Remark}

\author{Roshni Sahoo\\
  \texttt{r.sahoo@columbia.edu}}


\title{Density Ratio Estimation as Supervised Learning}

\date{Columbia University}

\begin{document}

\maketitle


\begin{section}{Introduction}

Density ratio estimation is a fundamental problem in causal inference and learning under distribution shift. The most common approach to density ratio estimation is probabilistic classification \citep{menon2016linking,sugiyama2008direct}. In probabilistic classification, we observe i.i.d. samples $X \sim P$ and $X \sim Q$. We provide pseudolabel $Y=1$ to samples drawn from $Q$ and pseudolabel $Y=0$ to samples drawn from $P$. Let $F$ denote the joint mixture distribution over $(X, Y)$ with these labels. Then, we fit a classifier $h(x) = \PP[]{Y=1 \mid X}$. By Bayes' rule, we can show that the density ratio $r(x) := \frac{dQ_{X}(x)}{dP_{X}(x)}$ is given by
\[ r(x) = \frac{h(x)}{1 - h(x)}.\]
The problem with this approach to density ratio estimation is that it can become unstable when classification probabilities approach $1$. Suppose that the classifier is estimated with some error $\hat{h}(x)$. Let $\delta(x) := \hat{h}(x) - h(x)$. We can compute the error in $r(x)$ that arises from the classifier error:
\begin{align*} 
|\hat{r}(x) - r(x)| &=  \Bigg| \frac{\hat{h}(x)}{1 - \hat{h}(x)} - \frac{h(x)}{1 - h(x)} \Bigg| \\
&= \Bigg| \frac{\delta(x) \cdot (1 + r(x))^{2}}{ 1 - \delta(x) \cdot ( 1 + r(x) ) } \Bigg|.
\end{align*}
Rather than estimate the classifier and then form the density ratio, we propose to directly estimate the density ratio by solving a loss minimization problem
\begin{equation}
\label{eq:direct_dre}
\min_{r: \mathcal{X} \rightarrow \mathbb{R}} \EE[F]{L_{\psi}(r(X), Y)}.
\end{equation}
Define $\psi:(0, \infty) \rightarrow\mathbb{R}$. Define $\phi(r) = \int_{0}^{r} t\psi'(t)dt.$ More broadly, we can consider the family of loss functions given below.
\begin{equation} 
\label{eq:loss}
L_{\psi}(r, y) = \phi(r) \cdot (1 - y) - \psi(r) \cdot y. \end{equation}

\begin{assumption}
\label{assumption:psi}
The function $\psi$ is twice-continuously differentiable, strictly increasing, concave, and has the property that $\psi'(r) + r\psi''(r) \geq 0.$
\end{assumption}

\begin{lemm}
Let $X_{i} \sim P$ denote i.i.d. samples from $P$ and $X_{j} \sim Q$ denote i.i.d. samples from $Q$. Pseudolabel the samples from $P$ by $Y_{i} = 0$ and pseudolabel the samples from $Q$ by $Y_{i}=1$ . The loss function $L_{\psi}$ given in \eqref{eq:loss} is convex in $r$. In addition, \eqref{eq:direct_dre} is strictly convex and has a unique minimizer at the true density ratio $r(x) = \frac{dQ_{X}(x)}{dP_{X}(x)}.$
\end{lemm}

\begin{proof}
We note that
\begin{align*}
L_{\psi}'(r, y) &= \phi'(r) \cdot (1 - y) - \psi'(r) \cdot y \\
&= r \psi'(r) \cdot (1 - y) - \psi'(r) \cdot y.
\end{align*}

In addition, we have that
\begin{align*}
L_{\psi}''(r, y) &= \phi''(r) \cdot (1 - y) - \psi''(r) \cdot y \\
&= (r\psi''(r) + \psi'(r)) \cdot (1 - y) - \psi''(r) \cdot y
\end{align*}

Since $\psi$ is concave, we have that $L_{\psi}''(r, 1) \geq 0$. Since $r\psi''(r) + \psi'(r) \geq 0$, we have that $L_{\psi}''(r, 0) \geq 0$. Thus, $L_{\psi}$ is convex in $r$.

We show that the unique population minimizer occurs at $r(x) = \frac{dQ_{X}(x)}{dP_{X}(x)}$. We have that the first-order condition of the population risk is given by 
\begin{align*}
\nabla_{r(x)} \EE[]{L_{\psi}(r(X), Y) \mid X=x} &= r(x) \cdot \psi'(r(x)) \cdot \PP[]{Y=0 \mid X=x} - \psi'(r(x)) \cdot \PP[]{Y=1 \mid X=x} \\
&= \psi'(r(x))  \cdot (r(x) \cdot \PP[]{Y=0 \mid X=x} - \PP[]{Y=1 \mid X=x}).
\end{align*}
Since $\psi$ is strictly increasing, $\psi'(r) > 0$, so the unique stationary point of the population risk is 
\begin{align*}
r(x) &= \frac{\PP[]{Y=1 \mid X=x}}{\PP[]{Y=0 \mid X=x}} =\frac{\PP[]{X=x \mid Y=1}}{\PP[]{X=x \mid Y= 0}} = \frac{dQ_{X}(x)}{dP_{X}(x)},\\
\end{align*}
where the last equalities follow by Bayes' rule.

\end{proof}

\end{section}


\bibliographystyle{plainnat} 
\bibliography{ref.bib}

\end{document}